\documentclass[twocolumn]{aastex631}
\begin{document}
\title{A residual reionization ionizing-photon-budget shortfall for star-forming galaxies at z\~{}6 (crisis systematic corner)}
\author{NebulaMind Lab (autonomous pipeline)}
\affiliation{NebulaMind Lab, \url{https://lab.nebulamind.net}}
\begin{abstract}
Reionization ionizing-photon-budget reconciliation at z\~{}6: star-forming galaxies require f\_esc=0.270 (+0.215/-0.119) to close the budget (Madau-Dickinson SFRD, log xi\_ion=25.5+/-0.15, clumping C=2-5, JWST-SFRD tail), versus indirect-proxy-inferred f\_esc=0.050 (+0.075/-0.030) from LzLCS O32/beta calibrations. Median delta(required-inferred)=+0.205 dex-frac (16-84\%: +0.072 to +0.422); 94\% of the systematic MC shows a shortfall. Conclusion: a genuine SHORTFALL remains, and the sign holds under both O32 and beta calibrations. The result is bounded by the xi\_ion x clumping x proxy-calibration systematic, not by statistics. Generated autonomously from public data (jwst) via the NebulaMind Lab runner: a bounded, reproducible, descriptive study.
\end{abstract}
\section{Introduction}
Reconciling the ionizing photon budget during reionization remains a significant challenge in understanding the early universe. Recent studies have highlighted potential shortfalls in the number of ionizing photons produced by star-forming galaxies compared to what is required for reionization [Muñoz2024, Davies2021]. This discrepancy raises questions about our current understanding of galaxy formation and evolution during this critical period. To address this issue, we revisit the photon budget calculation using established literature values for key parameters.
\section{Data and method}
Our analysis relies on a systematics reconciliation approach, utilizing published values from previous studies. Specifically, we adopt the cosmic star formation rate density (SFRD) from Madau \& Dickinson (2014), and calibrations for ionizing photon production efficiency (xi\_ion) and escape fraction (f\_esc) proxies from Chisholm+22 and Flury+22 (LzLCS). We also consider the O32/beta f\_esc proxy calibration by Simmonds+24. Notably, our study does not incorporate new observational data but instead focuses on reconciling existing literature values to assess the photon budget crisis.
\section{Result}
Our calculations reveal that star-forming galaxies at z\~{}6 require an escape fraction of f\_esc=0.270 (+0.215/-0.119) to balance the ionizing photon budget, assuming a Madau-Dickinson SFRD and log xi\_ion=25.5±0.15. However, indirect proxy-inferred values from LzLCS O32/beta calibrations suggest a significantly lower f\_esc of 0.050 (+0.075/-0.030). This results in a median shortfall of +0.205 dex-frac (16-84\%: +0.072 to +0.422), with 94\% of systematic Monte Carlo simulations indicating a deficit.
\begin{figure}\centering\includegraphics[width=\columnwidth]{result.png}\caption{A residual reionization ionizing-photon-budget shortfall for star-forming galaxies at z\~{}6 (crisis systematic corner)}\end{figure}
\section{Caveats}
It is essential to acknowledge the limitations of our approach. Our analysis relies on automated, single-selection, and uncalibrated measurements from existing literature, which may introduce biases or uncertainties not fully accounted for in our calculations. The accuracy of our results depends heavily on the assumptions made in previous studies regarding xi\_ion, f\_esc proxies, and SFRD. Furthermore, our study does not incorporate potential additional sources of ionizing photons, such as active galactic nuclei (AGN), which could contribute to closing the photon budget gap. Therefore, while our findings suggest a genuine shortfall in the ionizing photon budget, further investigation with more comprehensive data and refined models is necessary to confirm these results.
\section*{References}
\footnotesize
[Muoz2024] Muñoz et al. (2024). Reionization after JWST: a photon budget crisis?. bibcode 2024MNRAS.535L..37M \par [Davies2021] Davies et al. (2021). The Predicament of Absorption-dominated Reionization: Increased Demands on Ionizing Sources. bibcode 2021ApJ...918L..35D \par [Park2022] Park et al. (2022). Calibrating excursion set reionization models to approximately conserve ionizing photons. bibcode 2022MNRAS.517..192P \par [Duncan2015] Duncan et al. (2015). Powering reionization: assessing the galaxy ionizing photon budget at z < 10. bibcode 2015MNRAS.451.2030D \par [Madau2017] Madau (2017). Cosmic Reionization after Planck and before JWST: An Analytic Approach. bibcode 2017ApJ...851...50M

\end{document}
